\documentclass[acmsmall, screen]{acmart}

\settopmatter{printacmref=false}
\setcopyright{none}

\usepackage{eucal}

\usepackage{enumerate,xspace}
\usepackage[capitalise]{cleveref}

%% Bibliography style
\bibliographystyle{ACM-Reference-Format}
\citestyle{acmauthoryear}   %% For author/year citations

\begin{document}

\title{Constructing ODGI Graphs: Annotated Bibliography}
\author{Anshuman Mohan}
\affiliation{
\department{Bowers CIS}
\institution{Cornell University}
\city{Ithaca}
\country{USA}
}
\email{amohan@cs.cornell.edu}


\maketitle

\section{Introduction}

In the field of computational biology, there is much excitement about the
\emph{pangenome} graph. The traditional approach has long been to 
compare the given sample to a single reference genome, but this has a 
severe drawback: it is inherently overfitted to one reference.
In contrast, the pangenome allows for a many-to-many comparison of samples, 
thus allowing a much richer exploration of the similarities and differences
between organisms.

Once constructed, these graphs can be quieried to yield results that help 
identify aberrations in the human genome that may be the sources of diseases and
genetic proclivities. Used at its full potential, this technology could 
radically change medicial diagnoses and treatment. This has a 
direct tie to Sustainable Development Goal 3, ``Good Health and Well-Being'' 
\cite{sdg}.

Traditional genomic sequencing is famously expensive in terms of time and money, 
and the pangenome graph is even more so. The \emph{efficient} construction of 
pangenome graphs is an area of active development. 
This document surveys a few existing methods and summarizes their approaches. 
A more direct comparison, along with opportunities for a combination of these
approaches, will follow in future work.

\section{Constructing GFAs}

\subsection{Two Techniques}

The Graphical Fragment Assembly (GFA) format \cite{gfa} has 
emerged as the research standard for representing pangenome graphs.
Broadly, there are two techniques for the construction of these graphs. 

\cite{poa} 

\cite{genome_alignment}

\subsection{The Bleeding Edge}

\cite{minigraph}

\cite{progcactus}

\section{Variation graphs and ODGI}

\citet{variationgraphs}

Building atop of GFAs, \citet{odgi} introduce ODGI 
(Optimized Dynamic Genome/Graph Implementation),
which is (1) a way to represent pangenomes as variation graphs, 
and (2) a suite of commands that analyze these graphs efficiently. 
One of the novel contributions of the ODGI representation is that it does 
significant preprocessing on the graph and maintains a rich amount of 
information to hand. The set of nodes and the set of paths
are maintained at the top level, and further, every node maintains lists
of in- and out-edges, paths that traverse that node, and so on. 
This unlocks the potential for various kinds of parallelization
(nodewise, pathwise, or edgewise), depending on the command we wish to optimize. 
Note that ODGI's variation-graph-based representation is nonstandard; 
the ODGI project offers conversion methods to and from the GFA format.


\section{Monolithic Solutions}

\cite{pangolin}

\cite{panx}

\section{Efforts in other Domains}

\citet{drnets} or some other Carla paper, plus their genomic link.


\bibliography{gfa}

\end{document}